\documentclass[10pt,reqno]{amsart}
\usepackage[T1]{fontenc}
\usepackage{lmodern}
\usepackage{microtype}
\usepackage{amsmath,amssymb,mathtools}
\usepackage{booktabs}
\usepackage[colorlinks=true,linkcolor=blue,citecolor=blue,urlcolor=blue]{hyperref}
\newtheorem{theorem}{Theorem}[section]
\newtheorem{lemma}[theorem]{Lemma}
\newtheorem{proposition}[theorem]{Proposition}
\newtheorem{corollary}[theorem]{Corollary}
\theoremstyle{definition}
\newtheorem{definition}[theorem]{Definition}
\newtheorem{hypothesis}[theorem]{Hypothesis}
\newtheorem{remark}[theorem]{Remark}
\newtheorem{assessment}[theorem]{Assessment}
\newtheorem{problem}{Problem}
\numberwithin{equation}{section}
\newcommand{\Z}{\mathbb Z}
\newcommand{\N}{\mathbb N}
\newcommand{\Q}{\mathbb Q}
\newcommand{\R}{\mathbb R}
\newcommand{\HH}{\mathbb H}
\newcommand{\cL}{\mathcal L}
\newcommand{\cF}{\mathcal F}
\newcommand{\cT}{\mathcal T}
\DeclareMathOperator{\SL}{SL}
\DeclareMathOperator{\Br}{Br}
\DeclareMathOperator{\Nm}{N}
\newcommand{\status}[1]{\textup{\textsc{#1}}}
\newcommand{\leg}[2]{\left(\frac{#1}{#2}\right)}
\title[Automorphic and geometric reformulations of Erd\H{o}s--Straus]{Automorphic and geometric reformulations of the Erd\H{o}s--Straus equation, and a divisor-class barrier}
\author{Anonymous}
\date{Draft, October 2026; internal checks only, not externally refereed}
\emergencystretch=3em
\begin{document}
\begin{abstract}
For a prime $p$ we collect exact reformulations of the Erd\H{o}s--Straus
equation $4/p=1/x+1/y+1/z$ in automorphic and arithmetic-geometric
language. Type~I solutions are $\SL_2(\Z)$ matrices with a prescribed
entry, equivalently Heegner points of discriminant $-4d$ in a thin wedge,
whose discriminant-$(-4dp^2)$ forms are $T_p$-Hecke translates. Type~II
solutions are Heegner points of discriminant $-4dp$ lying on an exact
shell around a cuspidal geodesic. Further reformulations use orders in the
quaternion algebra $B_{p,\infty}$, the algebra $(-d,-p)_\Q$, a CM torus,
and an isotropic ternary form. The Bright--Loughran Brauer class is a
genus character; this is folklore and goes back to Yamamoto. All these
encodings reduce to one statement, which we prove: once two of the
Elsholtz--Tao coordinates are fixed, the solutions are either a finite
union of congruence classes of $p$, or the divisors of one explicit integer
$N_\theta(p)$ in one residue class. Every representative-independent
invariant of the automorphic objects is congruence or genus data. We then
record negative results, each with an exact scope. Fixed ternary embeddings
reduce to Mordell coverings. Any detector must work at main term
$p^{o(1)}$. Thickening a target is the same as averaging over $n$. Local,
genus and spinor solubility hold even at odd squares. Under Schinzel's
Hypothesis~H, every argument that certifies the equation through a fixed
finite set of slices fails infinitely often. We give the quantitative gaps
to known equidistribution theorems, list the escape hatches that remain
open, and report numerical evidence: $f_I(p)\ge1$ and $f_{II}(p)\ge1$ for
all primes $p\le10^8$. The Erd\H{o}s--Straus conjecture is \emph{not}
proved here. The negative statements concern the routes examined and known
technology. They are not impossibility theorems for the conjecture.
\end{abstract}
\maketitle
\setcounter{tocdepth}{1}
\tableofcontents

\section{Introduction}\label{sec:intro}
The Erd\H{o}s--Straus conjecture (ES) asserts that for every integer
$n\ge2$
\begin{equation}\label{eq:es}
 \frac4n=\frac1x+\frac1y+\frac1z,\qquad x,y,z\in\N,
\end{equation}
is soluble. It suffices to treat primes $n=p$, and by Mordell's
identities only primes in six classes modulo $840$ (the squares) remain
open. ES has been verified to very large bounds and holds for almost all
$n$ in a strong quantitative sense, but it is open for an individual large
prime.

\subsection*{Why automorphic and geometric input}
In earlier internal work of this campaign
(\cite{PO}, and Theorems~M and~C of \cite{PS}) we proved that
\emph{congruence-covering} arguments cannot settle ES pointwise. More
precisely, under Schinzel's Hypothesis~H, every correct bounded
witness-producing procedure built from ring operations, floor division,
sign tests and factorisations of the integers it meets fails for
infinitely many primes $p\equiv1\pmod{24}$ \cite[Thm~C]{PS}. This is the
procedure-level form of the classical odd-square obstructions (Yamamoto,
Mordell--Schinzel, Elsholtz--Tao \cite[Prop.~1.6]{ET}, Bright--Loughran
\cite[Cor.~1.4]{BL}). It leaves room for input that is \emph{not}
congruence data: spectral, automorphic or arithmetic-geometric
information about a fixed prime $p$. This note collects what such input
can see.

\subsection*{What this note contains}
\begin{enumerate}
\item \emph{Exact dictionaries} (\S\ref{sec:dict}, \status{proved}). Type~I
and Type~II solutions are rewritten as Heegner points, $\SL_2$ matrices,
Hecke translates, shells around cuspidal geodesics, orders in definite
quaternion algebras, divisors of an element of a CM field, and points on
an isotropic ternary quadric. Several of these are folklore. We mark
attributions explicitly; in particular, the identification of the
Bright--Loughran Brauer class with a genus character is \emph{not} new
(Yamamoto 1965, via \cite[App.~A]{BL}).
\item \emph{The divisor-class reduction} (\S\ref{sec:reduction},
\status{proved}). Each encoding is a ``slice'' of the Elsholtz--Tao
variety. Fixing two coordinates leaves either a congruence condition on
$p$ or the divisors of one integer $N_\theta(p)$ in one residue class
(Lemma~\ref{lem:slices}, Theorem~\ref{thm:reduction}).
\item \emph{Negative results with exact scope} (\S\ref{sec:negative}):
ternary embeddings, the scale constraint, thickening, local solubility at
odd squares, geometry, and Hypothesis-H-conditional emptiness of single
targets.
\item \emph{Quantitative gaps} to known equidistribution theorems
(\S\ref{sec:gaps}, \status{Assessment}), \emph{open problems}
(\S\ref{sec:open}), and \emph{numerical evidence} (\S\ref{sec:numerics}).
\end{enumerate}

\subsection*{What this note does not contain}
ES is not proved here, conditionally or unconditionally. We know of no
standard conjecture (GRH, GLH, Ramanujan, Selberg's eigenvalue conjecture,
twisted Linnik--Selberg) that is known to imply ES. The literature survey
behind this note \cite{LA} found no prior pointwise result for ES,
conditional or not. Every negative statement below concerns the specific
encodings and mechanisms analysed, or ``known technology''. None is a
theorem that ES cannot be proved by automorphic methods. Words like
``barrier'' are used in that mechanism-scoped sense; where they summarise
more than a proved statement, they are labelled \status{Assessment}.

\subsection*{Labels}
Each statement carries one of the labels \status{proved} (a complete proof
here, or in the cited internal file, relative to explicitly cited
literature), \status{conditional} (proved assuming a named hypothesis),
\status{certified} (an exact finite computation), \status{evidence}
(numerics), or \status{Assessment} (a judgement that is not a theorem).
Proofs of short items are given in full. For longer ones we give a sketch
and a pointer to the internal source file, which also records its internal
review. All internal files have been reviewed internally only; this note
has not been externally refereed.

\subsection*{Notation}
$p\ge5$ is prime, $\cL=\log p$, and $\chi_{D}=\leg{D}{\cdot}$ is the
Kronecker symbol. For a negative discriminant $D$, $\chi_D(-1)=-1$. We
write $[A,B,C]$ for $AX^2+BXY+CY^2$, and $f_I(p)$, $f_{II}(p)$ for the
numbers of Type~I and Type~II solutions in the normalisation of
\cite[\S2]{ET}.

%%CONTINUE

\begin{thebibliography}{99}
\bibitem[ET]{ET} C. Elsholtz and T. Tao, Counting the number of solutions to the Erd\H{o}s--Straus equation on unit fractions, \emph{J. Aust. Math. Soc.} \textbf{94} (2013), 50--105; arXiv:1107.1010.
\bibitem[BL]{BL} M. Bright and D. Loughran, Brauer--Manin obstruction for Erd\H{o}s--Straus surfaces, \emph{Bull. Lond. Math. Soc.} \textbf{52} (2020), 746--761; arXiv:1908.02526.
\bibitem[Ya]{Ya} K. Yamamoto, On the Diophantine equation $4/n=1/x+1/y+1/z$, \emph{Mem. Fac. Sci. Kyushu Univ. Ser. A} \textbf{19} (1965), 37--47. (Not accessed; cited through \cite[App.~A]{BL}.)
\bibitem[Mo]{Mo} L. J. Mordell, \emph{Diophantine equations}, Academic Press, 1969, pp.~287--290.
\bibitem[Sch]{Sch} A. Schinzel, On sums of three unit fractions with polynomial denominators, \emph{Funct. Approx. Comment. Math.} \textbf{28} (2000), 187--194.
\bibitem[Pi]{Pi} J. Pila, Density of integral and rational points on varieties, \emph{Ast\'erisque} \textbf{228} (1995), 183--187.
\bibitem[Du]{Du} W. Duke, Hyperbolic distribution problems and half-integral weight Maass forms, \emph{Invent. Math.} \textbf{92} (1988), 73--90.
\bibitem[DSP]{DSP} W. Duke and R. Schulze-Pillot, Representation of integers by positive ternary quadratic forms and equidistribution of lattice points on ellipsoids, \emph{Invent. Math.} \textbf{99} (1990), 49--57.
\bibitem[Yo]{Yo} M. P. Young, Weyl-type hybrid subconvexity bounds for twisted $L$-functions and Heegner points on shrinking sets, \emph{J. Eur. Math. Soc.} \textbf{19} (2017), 1545--1576; arXiv:1405.5457.
\bibitem[Hu]{Hu} P. Humphries, Equidistribution in shrinking sets and $L^4$-norm bounds for automorphic forms, \emph{Math. Ann.} \textbf{371} (2018), 1497--1543.
\bibitem[HR]{HR} P. Humphries and M. Radziwi\l\l, Optimal small scale equidistribution of lattice points on the sphere, Heegner points, and closed geodesics, \emph{Comm. Pure Appl. Math.} \textbf{75} (2022), 1936--1996; arXiv:1910.01360.
\bibitem[GZ]{GZ} B. H. Gross and D. B. Zagier, On singular moduli, \emph{J. Reine Angew. Math.} \textbf{355} (1985), 191--220.
\bibitem[GS]{GS} A. Ghosh and P. Sarnak, Integral points on Markoff type cubic surfaces, \emph{Invent. Math.} \textbf{229} (2022), 689--749; arXiv:1706.06712.
\bibitem[Le]{Le} H. W. Lenstra, Jr., Divisors in residue classes, \emph{Math. Comp.} \textbf{42} (1984), 331--340.
\bibitem[Ho]{Ho} C. Hooley, On the number of divisors of a quadratic polynomial, \emph{Acta Math.} \textbf{110} (1963), 97--114.
\bibitem[DFI]{DFI} W. Duke, J. Friedlander and H. Iwaniec, Equidistribution of roots of a quadratic congruence to prime moduli, \emph{Ann. of Math.} \textbf{141} (1995), 423--441.
\bibitem[FI]{FI} J. Friedlander and H. Iwaniec, \emph{Opera de cribro}, AMS Colloquium Publications \textbf{57}, 2010.
\bibitem[FI98]{FI98} J. Friedlander and H. Iwaniec, The polynomial $X^2+Y^4$ captures its primes, \emph{Ann. of Math.} \textbf{148} (1998), 945--1040.
\bibitem[HB]{HB} D. R. Heath-Brown, Primes represented by $x^3+2y^3$, \emph{Acta Math.} \textbf{186} (2001), 1--84.
\bibitem[Ma]{Ma} J. Maynard, Primes represented by incomplete norm forms, \emph{Forum Math. Pi} \textbf{8} (2020), e3.
\bibitem[ST]{ST} P. Sarnak and J. Tsimerman, On Linnik and Selberg's conjecture about sums of Kloosterman sums, in \emph{Algebra, Arithmetic, and Geometry}, Progr. Math. \textbf{270}, Birkh\"auser, 2009, 619--635.
\bibitem[PO]{PO} \emph{Pointwise obstructions for the Erd\H{o}s--Straus seed-component programme}, internal note, \texttt{paper/pointwise-obstruction.tex}.
\bibitem[PS]{PS} \emph{Formal genericity and pointwise $\Omega$-results}, internal note, \texttt{POINTWISE\_SIZE.md} (Theorems M and C), with internal review \texttt{reviews/pointwise-size-step1-review.md}.
\bibitem[TH]{TH} \emph{The Elsholtz--Tao Type~I sum is $O(N\log^2N)$}, internal note, \texttt{paper/es-typei-heegner-note.tex} (the bound (D)32b).
\bibitem[O123]{O123} \emph{Flip the average: fixed-$p$ spectral forms of Type~I}, internal note, \texttt{POINTWISE\_SPECTRAL.md}, review \texttt{reviews/pointwise-spectral-review.md}.
\bibitem[O124]{O124} \emph{Embedding ES into representations by quadratic forms}, internal note, \texttt{POINTWISE\_TERNARY.md}, review \texttt{reviews/pointwise-ternary-review.md}.
\bibitem[O126]{O126} \emph{Arithmetic geometry of the Erd\H{o}s--Straus surfaces}, internal note, \texttt{GEOMETRY\_ES.md}, review \texttt{reviews/geometry-es-review.md}.
\bibitem[O127]{O127} \emph{ES for large primes conditional on a standard conjecture?}, internal note, \texttt{CONDITIONAL\_ES.md}, review \texttt{reviews/conditional-es-review.md}.
\bibitem[O128]{O128} \emph{Automorphic structure of Type~II solutions for fixed $p$}, internal note, \texttt{POINTWISE\_TYPEII\_SPECTRAL.md}, reviews \texttt{reviews/pointwise-typeii-spectral-review.md}, \texttt{-review-r2.md}.
\bibitem[O130]{O130} \emph{Quaternionic / supersingular reformulation}, internal note, \texttt{QUATERNION\_ES.md}, review \texttt{reviews/quaternion-es-review.md}.
\bibitem[O131]{O131} \emph{Exact-formula search}, internal note, \texttt{EXACT\_FORMULA.md}, review \texttt{reviews/exact-formula-review.md}.
\bibitem[RT1]{RT1} \emph{Red team for the automorphic/geometric attack}, internal report, \texttt{reviews/redteam-automorphic.md}.
\bibitem[RT2]{RT2} \emph{Red team 2 and synthesis: the divisor-class reduction}, internal report, \texttt{reviews/redteam-automorphic-2.md}.
\bibitem[LA]{LA} \emph{Spectral and geometric tools for pointwise representation (literature survey)}, internal file, \texttt{LITERATURE\_AUTOMORPHIC.md}.
\end{thebibliography}
\end{document}
